\documentclass[12pt,a4paper]{article}
\usepackage[utf8]{inputenc}
\usepackage[T1]{fontenc}
\usepackage[brazil]{babel}
\usepackage{geometry}
\usepackage{setspace}
\geometry{margin=2.5cm}
\setstretch{1.15}
\setlength{\parindent}{0pt}
\setlength{\parskip}{0.35em}

\begin{document}

\begin{center}
\textbf{Ministério da Educação}\\
\textbf{Universidade Tecnológica Federal do Paraná -- Campus Pato Branco}\\
\textbf{Departamento Acadêmico de Informática}\\[0.5em]
\textbf{Disciplina: Redes de Computadores} \hfill \textbf{Prof. Dr. Eden Dosciatti}

\vspace{1.2em}
\Large\textbf{Atividade Avaliativa 5 -- UDP (Camada de Transporte)}
\end{center}

\vspace{1em}
\noindent\textbf{Nome:} Pedro Mariano dos Santos \hfill \textbf{RA:} 2562790

\vspace{1em}

\section*{Etapa 1: Captura real no Wireshark}

Nesta etapa, foi analisada uma consulta DNS (que usa UDP na porta 53) para observar os campos fixos do cabeçalho UDP.

\textbf{Pacote analisado:} foi analisado o par de pacotes DNS (consulta e resposta) gerado pelo comando \texttt{nslookup google.com}. No print detalhado do Wireshark, a resposta aparece com \texttt{Src Port 53} e \texttt{Dst Port 51937}, confirmando a inversão em relação ao pacote de consulta.

\subsection*{Campos do cabeçalho UDP}

\textbf{Source Port:} \texttt{51937} (na consulta DNS, porta efêmera do cliente).

\textbf{Destination Port:} \texttt{53}, porque é a porta padrão (well-known) reservada para o serviço DNS.

\textbf{Length:} \texttt{52 bytes}.

\textbf{Checksum:} \texttt{0x702a} (marcado como \textit{unverified} no Wireshark).

\textbf{Existe número de sequência ou confirmação (ACK) no UDP?}
Não. O UDP é um protocolo não orientado à conexão e não implementa confirmação de entrega, retransmissão nem controle de ordem na camada de transporte.

\section*{Etapa 2: Simulação no Packet Tracer}

A topologia foi montada com 1 PC, 1 switch e 1 servidor, com IPs estáticos:
\begin{itemize}
\item PC: \texttt{192.168.1.10/24}
\item Servidor: \texttt{192.168.1.100/24}
\item DNS no PC apontando para: \texttt{192.168.1.100}
\end{itemize}

No servidor DNS, foi criado o registro \texttt{www.udp.redes -> 1.1.1.1}. Em modo simulação, foi executado \texttt{nslookup www.udp.redes}.

\subsection*{Comportamento observado do UDP}

\textbf{Porta de origem (consulta):} \texttt{1025}.

\textbf{Porta de destino (consulta):} \texttt{53}.

\textbf{Inversão de portas na resposta:} sim, o servidor responde com origem \texttt{53} e destino na porta efêmera do cliente.

\textbf{Ausência de estabelecimento de conexão:} o primeiro pacote já carrega os dados da consulta, sem handshake prévio.

\textbf{Se a resposta se perder no caminho:} o servidor DNS não retransmite por conta própria, e o cliente aguarda até ocorrer timeout.

\section*{Respostas objetivas solicitadas}

\textbf{1) Qual a porta de destino usada pelo DNS?}\\
\texttt{53}.

\textbf{2) O UDP estabeleceu conexão antes de enviar os dados?}\\
Não, o envio ocorre sem fase de estabelecimento de conexão.

\textbf{3) Quantos campos existem no cabeçalho UDP visto no simulador?}\\
Quatro campos fixos: Source Port, Destination Port, Length e Checksum.

\section*{Conclusão}

A atividade mostrou, de forma prática, que o UDP é enxuto e rápido, mas não oferece garantias de entrega. Isso fica claro tanto na captura real do Wireshark quanto na simulação no Packet Tracer: o protocolo envia e recebe datagramas sem controle de conexão, sem confirmação e sem retransmissão automática.

\end{document}
